\chapter{Formulaire}

\noindent\textit{Les repères numériques et les rares formules du programme de 3ème,
regroupés par thème pour une révision rapide. Chaque valeur est expliquée et illustrée
dans le chapitre correspondant : ce formulaire n'est qu'un aide-mémoire, pas un
cours.}
\vspace{8pt}

\begin{multicols}{2}

\subsection*{\color{onbo}Chromosomes et caryotype}
\begin{itemize}
\item Caryotype humain normal : $46$ chromosomes $=$ $23$ paires ($22$ paires
d'autosomes $+$ $1$ paire de chromosomes sexuels XX ou XY)
\item Trisomie 21 : $47$ chromosomes (paire $21$ en triple exemplaire)
\end{itemize}

\subsection*{\color{onbo}Génétique (drépanocytose)}
\begin{itemize}
\item Génotypes : AA (sain), AS (porteur sain), SS (malade)
\item Croisement AS $\times$ AS : $\dfrac14$ AA, $\dfrac12$ AS, $\dfrac14$ SS
\item Système ABO : $3$ allèles (A, B, O) $\to$ $4$ groupes (A, B, AB, O)
\end{itemize}

\subsection*{\color{onbo}Mitose et fécondation}
\begin{itemize}
\item Mitose : $1$ cellule mère ($46$ chr.) $\to$ $2$ cellules filles identiques
($46$ chr. chacune)
\item Gamètes : $23$ chromosomes chacun
\item Fécondation : $23+23\to46$ chromosomes (cellule-œuf)
\end{itemize}

\subsection*{\color{onbo}Circulation sanguine}
\begin{itemize}
\item Artères : cœur $\to$ organes \quad Veines : organes $\to$ cœur
\item Infarctus du myocarde : obstruction d'une artère du \textbf{cœur}
\item AVC : obstruction/rupture d'un vaisseau du \textbf{cerveau}
\end{itemize}

\subsection*{\color{onbo}Respiration}
\begin{itemize}
\item Trajet de l'air : nez/bouche $\to$ pharynx $\to$ larynx $\to$ trachée $\to$
bronches $\to$ bronchioles $\to$ alvéoles
\item Fréquence respiratoire : $f=\dfrac Nt$ (mvt/min)
\item Débit ventilatoire : $D=V_c\times f$
\item Air inspiré $\approx21\,\%$ $\mathrm{O_2}$ ; air expiré $\approx16\,\%$
$\mathrm{O_2}$
\end{itemize}

\subsection*{\color{onbo}Excrétion urinaire}
\begin{itemize}
\item Trajet de l'urine : reins $\to$ uretères $\to$ vessie $\to$ urètre
\item Néphron : filtration $\to$ réabsorption
\item Pourcentage réabsorbé : $\%=\dfrac{V_f-V_u}{V_f}\times100$ ($V_f$ = volume
filtré, $V_u$ = volume d'urine excrétée)
\item Rendement rénal normal : $\approx99\,\%$ du liquide filtré est réabsorbé
\end{itemize}

\subsection*{\color{onbo}Thérapies}
\begin{itemize}
\item Sérothérapie : anticorps injectés, protection immédiate, courte durée
\item Vaccination : anticorps fabriqués par l'organisme, protection différée,
longue durée
\end{itemize}

\subsection*{\color{onbo}Paludisme et fièvre Ebola}
\begin{itemize}
\item Paludisme : parasite \textit{Plasmodium}, vecteur $=$ moustique anophèle
femelle
\item Fièvre Ebola : virus, transmission par contact direct (fluides corporels),
pas de vecteur
\end{itemize}

\subsection*{\color{onbo}Roches sédimentaires}
\begin{itemize}
\item Formation : altération $\to$ transport $\to$ sédimentation $\to$ diagenèse
(compaction $+$ cimentation)
\item Deux origines : détritique (débris de roches) ou biogène (restes
d'organismes)
\item Principe de superposition : strate la plus basse $=$ la plus ancienne
\end{itemize}

\end{multicols}
